\section{Materiais e métodos}\label{chap:metodos-final}

\subsection{Caracterização da pesquisa e fontes de evidência}

O trabalho caracteriza-se como pesquisa aplicada e tecnológica em Ciência da Computação, cujo principal resultado é a plataforma desenvolvida. Sua construção e avaliação envolveram consulta à literatura, definição das funcionalidades, implementação em módulos e verificações do funcionamento. A avaliação se concentra no software, não em um experimento agronômico. O estudo bibliográfico não seguiu um protocolo de revisão sistemática, e não se atribui ao projeto a adoção formal de \textit{Design Science Research}, pois esse procedimento não foi documentado.

Para descrever o que foi efetivamente desenvolvido, a revisão tomou o código-fonte como referência principal, em vez de assumir que todos os itens previstos estavam implementados. Foram examinados os arquivos que iniciam a aplicação, os módulos Python, o componente HTML/JavaScript, as regras de armazenamento e exportação, o conjunto de testes, as especificações e os documentos acadêmicos. A proposta original foi tratada como registro do planejamento. A relação entre objetivos, requisitos, funções, testes e seções do artigo está organizada em \path{docs/artigo_final/RASTREABILIDADE_TCC.md}, distinguindo o que foi implementado do que foi avaliado.

O desenvolvimento foi organizado em especificações de funcionalidades e tarefas realizadas em branches, isto é, ramificações do controle de versão, com revisão das alterações antes da integração. O histórico registra entregas de interface, armazenamento e exportação acompanhadas por testes. Na correção dos contratos de dados, novos testes foram executados antes das mudanças para reproduzir as divergências e, depois, verificar se haviam sido corrigidas. Esse procedimento localizado não permite afirmar que todo o projeto seguiu desenvolvimento orientado a testes (TDD, \textit{Test-Driven Development}). O controle das dependências e a divisão de responsabilidades entre módulos apoiaram a evolução incremental; não foi identificada configuração de integração contínua no repositório.

\subsection{Materiais e ambiente}

A solução foi implementada em Python. O Streamlit coordena a interface e o estado da sessão, enquanto Folium/Leaflet apresentam os mapas e o Earth Engine executa as operações raster. Após a transferência dos resultados, NumPy e pandas organizam os dados locais e Plotly produz os gráficos. Para disponibilizar esses resultados fora da interface, ReportLab, Pillow e Matplotlib apoiam os relatórios estáticos, e openpyxl gera planilhas. A biblioteca geemap permanece em operações e mapas auxiliares, mas não é responsável por toda a interface atual.

O arquivo \texttt{requirements.txt} fixa dependências diretas e transitivas. A execução da suíte registrada nas evidências utilizou Linux x86\_64, Python 3.14.4, pytest 9.0.3, pytest-cov 7.1.0 e coverage 7.14.0. A identificação dessas versões situa o resultado dos testes, mas não demonstra reprodução em todos os sistemas operacionais ou na hospedagem. Os registros de preparação do Streamlit Community Cloud descrevem separadamente a resolução de dependências para Python 3.12.

\subsection{Entradas e requisitos}

As entradas principais são um polígono, duas datas e a escolha de NDVI, NDWI, NDMI ou todos os índices. A busca textual serve para navegação, não para inferir automaticamente limites de uma propriedade. A análise depende de geometria válida e intervalo cuja data final não anteceda a inicial. Salvar uma localização não executa análise nem associa automaticamente uma cultura.

O sistema permite selecionar a área, consultar imagens Sentinel-2, aplicar filtros, calcular índices e apresentar mapas, séries, sinalizações de anomalias e informações meteorológicas. Também permite nomear e reutilizar áreas e exportar resultados. Quanto às características de qualidade, a divisão em módulos pode ser examinada no código e a cobertura dos testes pode ser medida. Já a simplicidade da interface, o tempo de resposta inferior a 30 s, o baixo consumo e a redução de custos não foram demonstrados por avaliações específicas. Na execução local, a própria máquina precisa executar o servidor Streamlit, sem exigir um servidor dedicado separado.

\subsection{Aquisição de imagens e qualidade}

O módulo \texttt{earth\_engine.py} consulta a coleção \texttt{COPERNICUS/S2\_SR\_HARMONIZED} e mantém as imagens que possuem interseção com a área selecionada. Na interface, as datas inicial e final fazem parte do período solicitado, em Tempo Universal Coordenado (UTC). O Earth Engine, porém, não inclui o instante final de \texttt{ee.Filter.date}. Por isso, a aplicação envia o início do dia seguinte como limite: $[d_{\mathrm{inicial}},d_{\mathrm{final}}+1\ \mathrm{dia})$ \citep{GOOGLE_EE_DATE}. Assim, selecionar a mesma data nos dois campos corresponde a um dia completo. Uma data final anterior à inicial é rejeitada.

Antes de calcular os índices, o sistema filtra as imagens pela informação de cobertura de nuvens do produto. A condição padrão é estrita: \texttt{CLOUD\_COVERAGE\_ASSESSMENT < 20}, portanto o valor deve ser inferior a 20. Essa informação descreve o arquivo do produto que contém o grânulo, não uma medição de nuvens feita dentro da área selecionada. Todas as imagens aprovadas nos filtros seguem para o processamento; não é escolhida apenas a mais recente ou a menos nublada.

Depois desse filtro, a banda de qualidade QA60 é usada para excluir posições marcadas como nuvens opacas ou cirrus. A máscara mantém somente as posições em que os bits 10 e 11 dessa banda são zero. O conjunto selecionado contém B2, B3, B4, B8, B8A, B11 e QA60, embora B2 e B3 não entrem nas três equações dos índices.

Esse procedimento não inclui uma máscara adicional específica para sombras, neve ou probabilidade de nuvens. Há também uma restrição histórica: o catálogo informa que QA60 esteve mascarada entre 25/01/2022 e 28/02/2024, inclusive, e descreve sua posterior reconstrução a partir de classificações auxiliares \citep{GOOGLE_S2_SR}. Assim, encontrar imagens na coleção não garante que restarão dados úteis após aplicar a máscara em qualquer período.

Os índices das \cref{eq:ndvi-final,eq:ndwi-final,eq:ndmi-final} são calculados por \texttt{normalizedDifference}, renomeados e associados a \texttt{system:time\_start} e \texttt{system:index}. A série preserva data, valor e identificador de aquisição quando disponível, permitindo distinguir observações do mesmo dia.

\subsection{Produtos espaciais e temporais}\label{sec:agregacoes-final}

O processamento produz um mapa que combina o período e uma série que mantém os resultados de cada aquisição. Para formar o mapa, o sistema calcula primeiro o índice em cada imagem e depois combina seus valores por mediana em cada posição. Representando por $I_k(p,t)$ o índice $k$ na posição $p$ e aquisição $t$, após as máscaras, a composição é:
\begin{equation}\label{eq:composto-final}
 C_k(p)=\operatorname{mediana}_{t\in T_p} I_k(p,t),
\end{equation}
em que $T_p$ reúne as observações válidas naquela posição. A ordem é importante: calcula-se a mediana dos índices já obtidos, não o índice a partir de uma mediana prévia das bandas.

Para apresentar um único valor junto ao mapa, o sistema calcula a média espacial da composição $C_k$. Para construir a série temporal, calcula a média espacial separadamente em cada imagem. As duas operações usam o redutor de média do Earth Engine, mas recebem imagens diferentes. Elas podem ser representadas por:
\begin{align}
 M_k&=\frac{\sum_{p\in P_C}w_{p,C}C_k(p)}
                {\sum_{p\in P_C}w_{p,C}},\label{eq:media-composto}\\
 S_k(t)&=\frac{\sum_{p\in P_t}w_{p,t}I_k(p,t)}
                  {\sum_{p\in P_t}w_{p,t}}.\label{eq:serie-final}
\end{align}
Nessas expressões, $P_C$ e $P_t$ são os conjuntos de posições válidas, e $w$ indica quanto cada posição contribui para a média. A soma dos pesos, no denominador, deve ser positiva. No redutor ponderado padrão, o peso é o menor valor entre a máscara e a fração aproximada do pixel coberta pela área, não o produto dos dois. Essa fração é representada internamente em oito bits; interseções inferiores a aproximadamente $1/256$ podem receber peso zero \citep{GOOGLE_EE_REGION}.

A aplicação não acrescenta outra ponderação ao resultado. A chamada é \texttt{ee.Reducer.mean()}, com escala de processamento \texttt{scale=10} e limite \texttt{maxPixels=1e9}, sem definição explícita de \texttt{crs} nessa etapa \citep{GOOGLE_EE_REDUCE}.

Cada resultado responde a uma pergunta diferente. $C_k(p)$ informa o valor de uma posição no mapa composto; $M_k$ resume esse mapa em uma média da área; e $S_k(t)$ informa a média da área em uma aquisição. Esses valores não devem ser usados como se fossem equivalentes. Também é necessário considerar a grade espacial: o sistema de referência de coordenadas (CRS) não é fixado explicitamente nas reduções anteriores. A grade em projeção Universal Transversa de Mercator (UTM), com escala de 10 m, é definida depois, para a consulta interativa.

Ao organizar a série em \texttt{time\_series.py}, o sistema omite os registros sem valor retornado pelo serviço, converte os demais em números e os ordena pela data no formato ISO. Não estima valores entre as observações, não preenche lacunas e não faz agrupamento mensal ou composição diária. Imagens diferentes de uma mesma data podem gerar registros distintos, pois datas repetidas não são eliminadas. A contagem de imagens da coleção é feita antes da máscara; por isso, pode ser maior que a quantidade de pontos válidos na série.

\begin{figure*}[htbp]
\centering
\resizebox{0.96\textwidth}{!}{\begin{tikzpicture}[
 box/.style={draw,rounded corners,align=center,text width=36mm,
 minimum height=13mm,font=\small,fill=green!5},
 link/.style={-{Stealth[length=2mm]},thick}]
\node[box] (input) at (0,0) {Polígono, período\\e índices solicitados};
\node[box] (query) at (4.8,0) {Coleção Sentinel-2\\interseção, datas\\e nuvens do produto};
\node[box] (indices) at (9.6,0) {Máscara QA60\\diferença normalizada\\em cada aquisição};
\node[box] (spatial) at (14.4,1.6) {Mediana temporal\\por pixel\\$C_k(p)$};
\node[box] (mean) at (19.2,1.6) {Média espacial\\da composição\\$M_k$};
\node[box] (series) at (14.4,-1.6) {Média espacial\\por imagem\\$S_k(t)$};
\node[box] (anomaly) at (19.2,-1.6) {Série ordenada\\tendência e desvios\\sem diagnóstico};
\node[box,fill=blue!5] (climate) at (4.8,-3.7) {Centroide\\NASA POWER\\tabela e resumo};
\node[box] (output) at (14.4,-3.7) {Resultados por índice\\mapa, gráfico e snapshot\\para exportação};
\draw[link] (input) -- (query);
\draw[link] (query) -- (indices);
\draw[link] (indices) -- (spatial);
\draw[link] (spatial) -- (mean);
\draw[link] (indices) -- (series);
\draw[link] (series) -- (anomaly);
\draw[link] (input.south) |- (climate.west);
\draw[link] (climate) -- node[above,font=\scriptsize] {contexto compartilhado} (output);
\draw[link] (anomaly.south) |- (output.east);
\draw[link] (mean.east) -- ++(0.8,0) |- (output.south east);
\draw[link] (spatial.south west) -- ++(-0.6,-1.4) -- (output.north west);
\end{tikzpicture}
}
\caption{Derivação dos produtos espaciais e temporais a partir da coleção filtrada. Fonte: elaborado pelo autor a partir de \texttt{pipeline.py} e \texttt{time\_series.py}.}
\label{fig:fluxo-analise}
\end{figure*}

\subsection{Regra de detecção e tendência}\label{sec:anomalias-final}

Para sinalizar quedas na série, o sistema compara cada valor com os três registros imediatamente anteriores. Primeiro, calcula a média e o desvio padrão desses registros; depois, verifica quanto o valor atual se afasta dessa referência. Esse afastamento é expresso pelo escore $z_i$. A janela considera três observações, não três dias, pois o intervalo entre imagens pode variar. Nas equações seguintes, $x_i$ representa o valor na posição $i$ da série, com contagem iniciada em zero. O cálculo pode ser aplicado a partir da quarta observação, quando $i\geq3$:
\begin{align}
 \mu_i&=\frac{1}{3}\sum_{j=i-3}^{i-1}x_j,\\
 \sigma_i&=\sqrt{\frac{1}{3}\sum_{j=i-3}^{i-1}(x_j-\mu_i)^2},\\
 z_i&=\frac{x_i-\mu_i}{\sigma_i},\quad \sigma_i>0.\label{eq:escore-final}
\end{align}
O algoritmo marca anomalia somente quando $z_i<-1$: o valor atual deve estar mais de um desvio padrão abaixo da média dos três anteriores. Os três primeiros registros não têm uma janela completa e não são sinalizados. O mesmo ocorre quando $\sigma_i=0$, pois não é possível dividir o desvio por zero. Nesses casos, a saída armazena escore zero por convenção do algoritmo, não como resultado de um escore calculado.

O desvio padrão usa o divisor $N$, correspondente à opção \texttt{ddof=0} do NumPy, em vez de $N-1$ da fórmula amostral \citep{NUMPY_STD}. A denominação populacional dessa fórmula não significa que os três valores representem todas as condições agrícolas possíveis. A média usa apenas os registros anteriores e exclui o valor atual.

A função \texttt{detect\_anomalies} pressupõe entradas numéricas compatíveis, sem tratamento adicional específico de NaN, marcador de valor não numérico. Valores já sinalizados continuam participando das janelas seguintes. A janela de três registros e o limiar $-1$ são escolhas heurísticas, isto é, regras particulares da aplicação sem calibração experimental documentada. Elas não modelam sazonalidade nem são ajustadas por cultura.

Para resumir a tendência, o sistema ajusta uma reta à sequência de valores. A função \texttt{np.polyfit} calcula $x_i\simeq a\,i+b$ pelo método dos mínimos quadrados, minimizando $\sum_i[x_i-(a\,i+b)]^2$ \citep{HYNDMAN2021,NUMPY_POLYFIT}. O coeficiente $a$ é a inclinação, que indica a variação do índice por posição na sequência; $b$ é o intercepto, o valor da reta na posição zero.

A tendência é classificada como crescente se $a>0{,}01$, decrescente se $a<-0{,}01$ e estável nos demais casos. Esses limites também são heurísticos, sem calibração documentada. Como o ajuste usa a ordem das observações e não os dias transcorridos, a inclinação não é uma taxa diária. Quando há menos de dois pontos, a função retorna o rótulo estável por convenção, embora os dados sejam insuficientes para estimar uma inclinação.

\begin{figure*}[htbp]
\centering
\resizebox{0.92\textwidth}{!}{\begin{tikzpicture}[
 node distance=9mm and 12mm,
 box/.style={draw,rounded corners,align=center,text width=40mm,
 minimum height=12mm,font=\small,fill=green!5},
 decision/.style={draw,diamond,aspect=2,align=center,font=\small,inner sep=1mm},
 link/.style={-{Stealth[length=2mm]},thick}]
\node[box] (input) {Série ordenada\\valores $x_i$};
\node[decision,right=of input] (initial) {$i<3$?};
\node[box,right=of initial] (window) {Três valores anteriores\\$\mu_i$ e $\sigma_i$ populacional};
\node[decision,right=of window] (zero) {$\sigma_i=0$?};
\node[box,below=of window] (score) {$z_i=(x_i-\mu_i)/\sigma_i$\\sinaliza se $z_i<-1$};
\node[box,below=of initial] (false) {Sem sinalização\\escore inicial zero};
\node[box,below=of score] (output) {Registro com data, valor,\\\texttt{anomaly} e \texttt{z\_score}\\avança para próximo ponto};
\draw[link] (input)--(initial);
\draw[link] (initial)--node[above,font=\scriptsize]{não}(window);
\draw[link] (initial)--node[left,font=\scriptsize]{sim}(false);
\draw[link] (window)--(zero);
\draw[link] (zero.south)|-node[above,font=\scriptsize,pos=.6]{não}(score.east);
\draw[link] (zero.north)--++(0,0.8)-|node[above,font=\scriptsize,pos=.5]{sim}(false.west);
\draw[link] (false.south)|-(output.west);
\draw[link] (score)--(output);
\end{tikzpicture}
}
\caption{Heurística unilateral de sinalização de desvios negativos. Fonte: elaborado pelo autor a partir de \texttt{anomalies.py}.}
\label{fig:fluxo-anomalias}
\end{figure*}

\subsection{Consulta e agregação meteorológica}\label{sec:clima-final}

O contexto meteorológico é obtido por uma consulta diária da NASA POWER nas coordenadas do centroide calculado pelo Earth Engine. A requisição usa o serviço \texttt{/api/temporal/daily/point}, a comunidade agrícola \texttt{community=ag}, datas no formato \texttt{YYYYMMDD}, resposta JSON e tempo limite de 30 s. São solicitados os parâmetros \texttt{PRECTOTCORR}, \texttt{T2M} e \texttt{ALLSKY\_SFC\_SW\_DWN}. Embora o código contenha nomes internos para \texttt{T2M\_MAX} e \texttt{T2M\_MIN}, o fluxo padrão não consulta temperaturas máximas e mínimas.

No catálogo diário da comunidade agrícola, \texttt{PRECTOTCORR} corresponde à precipitação corrigida, em mm/dia; \texttt{T2M}, à temperatura média do ar a 2 m, em °C; e \texttt{ALLSKY\_SFC\_SW\_DWN}, à energia solar de ondas curtas incidente em superfície horizontal sob todas as condições de céu, em MJ/m²/dia \citep{NASA_POWER_PARAMETERS}. Para a integração efetiva, prevalece a unidade declarada em cada resposta, especialmente na radiação.

A resposta é organizada em uma tabela ordenada por data. O código $-999$, usado para indicar ausência de dado, não é tratado como uma medição. Para resumir o período, o sistema soma a precipitação e calcula as médias de temperatura e radiação, usando separadamente as observações válidas de cada variável. Esses conjuntos de dias são representados por $D_P$, $D_T$ e $D_R$:
\begin{align}
 P_{\mathrm{total}}&=\sum_{d\in D_P} P_d,\\
 \overline{T}&=\frac{1}{|D_T|}\sum_{d\in D_T}T_d,\\
 \overline{R}&=\frac{1}{|D_R|}\sum_{d\in D_R} f(u)R_d.\label{eq:radiacao-codigo}
\end{align}
Na \cref{eq:radiacao-codigo}, $f(u)=1$ para MJ/m²/dia e $f(u)=3{,}6$ para kWh/m²/dia, conforme a unidade declarada nos metadados. A conversão decorre de $1\ \mathrm{kWh}=1000\ \mathrm{J/s}\times3600\ \mathrm{s}=3{,}6\ \mathrm{MJ}$, pelas definições de watt e prefixos do SI \citep{BIPM_SI}. Unidade ausente ou desconhecida torna a radiação indisponível e gera aviso, sem inferência pela magnitude. Valores não finitos, booleanos, inválidos e a sentinela $-999$ são descartados. A coluna normalizada em MJ/m²/dia é compartilhada pelo resumo e pelas exportações.

A soma da precipitação disponível não assegura cobertura integral do intervalo caso haja lacunas. A requisição define \texttt{time-standard=UTC}, mantendo as duas datas inclusivas em correspondência ao período espectral \citep{NASA_POWER_DAILY}. A consulta é compartilhada entre os índices da mesma análise. A apresentação conjunta fornece contexto, sem calcular correlação ou atribuir causalidade.

\subsection{Procedimentos de avaliação técnica}\label{sec:avaliacao-final}

A avaliação procurou verificar tanto regras locais quanto a comunicação entre as partes da aplicação. Testes unitários examinaram regras específicas; testes de integração verificaram a troca de dados entre componentes; e testes de regressão verificaram a manutenção de comportamentos já corrigidos. Essa distinção de escopos, assim como a diferença entre executar código e verificar seu comportamento, é discutida por \citet{BENDER2020}. A referência não implica adoção dos processos internos do Google.

A interação com serviços externos e a navegação em navegador foram registradas separadamente. Essa separação é necessária porque substituir um serviço durante um teste permite controlar sua resposta, mas não demonstra que o fluxo completo funciona com o serviço real.

Foi reexecutada a suíte existente com tracer C e cobertura de ramificações. Os testes espectrais verificam seleção de bandas, filtros e chamadas usando substitutos da API, não a aritmética remota sobre imagens reais. Os testes de persistência usam arquivos temporários; os de exportação inspecionam formatos e valores controlados. Alguns testes Streamlit usam widgets reais com serviços substituídos. A cobertura Python não mede execução JavaScript no navegador.

Uma requisição pública curta à NASA POWER verificou metadados e datas em uma coordenada técnica, sem identificação de cultura. A resposta histórica está em tempo solar local (LST, \textit{Local Solar Time}); foi preservada para reproduzir a divergência de unidade, não como evidência de uma nova consulta UTC. Seus valores foram reutilizados em testes de regressão de resumo e exportação. O componente JavaScript recebeu testes independentes de legenda, consulta e descompressão, e matrizes controladas permitiram medir o empacotamento de dados. Essas medições não foram usadas para inferir desempenho em produção.

\subsection{Avaliação com potenciais usuários}\label{sec:pesquisa-usuarios}

A avaliação com usuários buscou conhecer as tarefas que os respondentes disseram ter realizado, sua experiência de uso e a utilidade atribuída aos recursos. Trata-se de uma avaliação exploratória e descritiva, com análise quantitativa e qualitativa, não de uma verificação da eficácia agronômica da plataforma. Conforme informação do autor, o formulário \url{https://forms.gle/VXqtdCgH9us774VE6} foi divulgado por WhatsApp e LinkedIn para profissionais e estudantes de agronomia, produção rural, tecnologia e áreas afins.

Como cada pessoa decidiu participar voluntariamente, a coleta caracteriza amostragem não probabilística por adesão, sujeita à autoseleção \citep{AAPOR2013}. Não há um cadastro do público convidado que permita calcular a taxa de resposta ou afirmar representatividade populacional.

A apresentação solicita exploração prévia da plataforma, estima de 5 a 8 minutos para responder e orienta evitar informações pessoais. A planilha fornecida pelo autor contém uma aba de respostas com 50 colunas, incluindo o carimbo, a triagem e os itens de matrizes expandidos, e outra aba vazia. Foram identificados três itens de perfil, nove atividades, nove afirmações sobre uso, onze recursos avaliados quanto à utilidade, oito afirmações sobre compreensão e aplicabilidade, cinco campos discursivos e uma nota geral de 0 a 10. O conhecimento declarado refere-se à agricultura de precisão, não a uma medida direta de domínio de sensoriamento remoto; a experiência prévia e sua frequência são captadas em um mesmo item.

Antes dos cálculos, as respostas foram conferidas quanto à elegibilidade, ao preenchimento e à duplicidade. Os 31 registros indicam ``Sim'' na pergunta conjunta sobre maioridade, leitura e consentimento, e na declaração de exploração de ao menos um recurso. Não há duplicatas exatas, nem respostas integralmente iguais ao desconsiderar o carimbo. Não foram coletados identificadores individuais, e a unidade da análise descritiva é a submissão elegível ao questionário.

O conjunto analisado abrange respostas registradas entre 06 e 15 de novembro de 2026.

Nas tarefas, as categorias são ``Consegui sem ajuda'', ``Consegui com ajuda'', ``Tentei, mas não consegui'' e ``Não tentei''. A conclusão entre tentativas é a razão entre as duas primeiras e o total das três primeiras. O questionário agrega séries temporais e possíveis anomalias em uma atividade e reúne a exportação de relatório ou dados em outra; não se estimam taxas separadas para recursos que o instrumento não separou. Trata-se de autorrelato, sem observação controlada, cronometração ou inspeção dos logs de cada utilização.

As perguntas de concordância usam cinco categorias, de discordo totalmente a concordo totalmente. Nas perguntas de utilidade, as categorias vão de nada útil a muito útil; ``Moderadamente útil'' ocupa a posição intermediária, mas não significa neutralidade. Para respeitar essa ordenação, a análise prioriza a contagem e o percentual das respostas em cada categoria, a mediana e os quartis, conforme as cautelas discutidas por \citet{SULLIVAN2013}. Os agrupamentos 4--5, 3 e 1--2 ajudam a resumir os resultados sem substituir as distribuições completas.

A base de cálculo depende da pergunta. ``Não consegui avaliar'', ausências e valores inválidos são contados separadamente e não entram no total de respostas usado para os percentuais ordinais. Nas tarefas, ``Não tentei'' permanece na distribuição geral, mas fica fora do cálculo de conclusão entre tentativas. Não há preenchimento estimado de respostas ausentes.

Os quartis usam interpolação linear, e o intervalo interquartil é $Q_3-Q_1$. Médias e desvios padrão amostrais dos códigos ordinais são disponibilizados somente como descrições complementares; a nota geral recebe também moda, extremos e distribuição. Não se construíram escores agregados de usabilidade ou aceitação, nem se denominou o instrumento SUS, TAM, UTAUT ou NPS. Seis cruzamentos descritivos usam agrupamentos amplos, com supressão de grupos com menos de cinco respostas válidas, sem testes de hipótese ou conclusões causais. Isso reduz a exposição de pequenos grupos, mas não certifica anonimato.

As respostas abertas foram lidas integralmente e organizadas por temas, em diálogo com a referência de análise temática de \citet{BRAUN2006}. Foram identificados trechos com significado relevante, atribuídas categorias conforme seu conteúdo e revistos os agrupamentos. Uma resposta podia receber mais de uma categoria. Cada registro, porém, conta apenas uma vez em cada tema, mesmo quando o menciona em perguntas diferentes.

Foram codificadas 142 células preenchidas, e treze comentários opcionais estavam ausentes. O campo sobre impedimento de acesso estava vazio nos 31 registros; essa ausência foi considerada compatível com a triagem do formulário, não uma medição de dificuldade igual a zero. A codificação teve assistência de IA, e seus critérios e trechos literais sem identificação permanecem disponíveis para validação do autor. Não houve codificação independente por dois avaliadores nem medição de concordância entre eles.

A coleta utiliza a versão pública anterior às correções de radiação, legenda, clima multiíndice, intervalos, invalidação e IDs de aquisição. Por decisão do autor, essa versão foi mantida durante a coleta, sem implantação das alterações locais. As percepções dos respondentes não são atribuídas à versão aprimorada, consolidada no commit \texttt{28759bf}. A documentação da análise distingue essas versões; o commit exato da versão pública permanece pendente. A confirmação dos procedimentos éticos institucionais cabe ao autor e ao orientador; não se atribui aprovação ou dispensa formal sem documento.
